\documentclass{article}
% ===== Preamble (XeLaTeX) =====
\usepackage{fontspec}
\usepackage{polyglossia}
\setmainlanguage{russian}
\setotherlanguage{english}
\defaultfontfeatures{Ligatures=TeX}

\setmainfont{DejaVuSerif.ttf}[
  BoldFont       = DejaVuSerif-Bold.ttf,
  ItalicFont     = DejaVuSerif-Italic.ttf,
  BoldItalicFont = DejaVuSerif-BoldItalic.ttf,
]
\setmonofont{DejaVuSansMono.ttf}[
  BoldFont       = DejaVuSansMono-Bold.ttf,
  ItalicFont     = DejaVuSansMono-Oblique.ttf,
  BoldItalicFont = DejaVuSansMono-BoldOblique.ttf,
]

% поля уменьшены примерно вдвое от стандартных article
\usepackage[a4paper,margin=1.9cm]{geometry}

\usepackage{amsmath,amssymb,amsfonts}
\usepackage{graphicx}
\usepackage{booktabs}
\usepackage{longtable}
\usepackage{array}
\usepackage{tabularx}
\usepackage{hyperref}
\usepackage{url}

\begin{document}

\title{DSIC --- Результаты эмпирической проверки\\(Часть I + Часть II + Часть III)}
\date{}
\maketitle

\noindent\textbf{Модель:} Dual Scale Inversion Cosmology (DSIC)\\
\textbf{Автор:} Азамат Засеев\\
\textbf{Скрипт:} \texttt{dsic\_test.py} (единый, самодостаточный; блоки \texttt{[0]}--\texttt{[11]})\\
\textbf{Сравнение:} DSIC против плоской $\Lambda$CDM на идентичных данных, при равном числе свободных параметров.

\begin{quote}
\textbf{Краткий вывод.} Внутренние тождества ядра проходят машинную проверку с точностью до $10^{-8}$. На геометрии расширения поздней Вселенной ($z < 2.3$) DSIC с единственным параметром $\mu_0 = 0.7600 \pm 0.0091$ статистически неотличима от $\Lambda$CDM на сверхновых и слегка предпочтительнее на совместном фите SN+BAO; $\Lambda$CDM-ветка конвейера воспроизводит опубликованные значения Pantheon+SH0ES --- валидация пайплайна. $Om(z)$-диагностика на текущих данных модели \textbf{не разрешает} (знак $\Delta\chi^2$ зависит от нормировки $H_0$). Второй этаж (порог отлипания $\mu_d$) точно замыкает расстояние до реликта, не затрагивая позднюю Вселенную тождественно; статус --- калибровка при форме, зафиксированной постулатом P6. Рост структур ($f\sigma_8$) на фоне DSIC в GR-пределе описывает компиляцию Gold-2017 не хуже $\Lambda$CDM ($\chi^2/\text{dof} \approx 0.7$); данные вырождены по $(\Omega_{m0}, \sigma_8)$, поэтому докладывается инвариант $S_8 = 0.761 \pm 0.013$ --- статус проверки на выживание, не вывода из ядра. Самая жёсткая зона риска, предъявленная числом, --- сжатое космическое время при $z = 5\text{--}12$ (при $z = 7.5$ доступно 0.42 Гыр против 0.71 у $\Lambda$CDM). Результаты воспроизводимы: скрипт сам скачивает публичные каталоги и считает $\chi^2$ с полной ковариацией.
\end{quote}
\clearpage
\section{Сводка результатов}

\begin{center}
\scriptsize
\setlength{\tabcolsep}{4pt}
\renewcommand{\arraystretch}{1.25}
\begin{tabularx}{\textwidth}{@{}>{\raggedright\arraybackslash}p{2.5cm} >{\raggedright\arraybackslash}X >{\raggedright\arraybackslash}p{1.9cm} >{\raggedright\arraybackslash}p{2.2cm} >{\raggedright\arraybackslash}p{2.3cm} >{\raggedright\arraybackslash}X@{}}
\toprule
\textbf{Проба} & \textbf{Данные / блок} & \textbf{Метрика} & \textbf{DSIC} & \textbf{$\Lambda$CDM} & \textbf{Вердикт} \\
\midrule
Связность ядра & тождества, \texttt{[0b]} & макс.\ откл. & $\le 5.7\cdot10^{-8}$ & --- & все тождества выполняются \\
Сверхновые (якорь + $M$ явно) & Pantheon+SH0ES (1657 SN, 77 калибр.), \texttt{[1]} & $\chi^2/\text{dof}$ & 0.8783 & 0.8783 & неотличимо ($\Delta\chi^2 = +0.05$) \\
Погрешности параметров & гессиан $\chi^2$, \texttt{[1]} & $1\sigma$ & $\mu_0 = 0.7600 \pm 0.0091$ & $\Omega_m = 0.3308 \pm 0.0180$ & параметр измерен ($\sim$1.2\%) \\
Устойчивость к ошибкам & STAT-ONLY, \texttt{[2]} & $\Delta\chi^2$ & --- & --- & устойчиво ($\Delta\chi^2 = -0.39$) \\
SN + BAO совместно & + DESI DR2 (6 бинов), \texttt{[3]} & $\Delta\chi^2$(тотал) & --- & --- & DSIC лучше на $-5.0$ (слабое) \\
BAO $D_M$--$D_H$ & DESI DR2 ($12{\times}12$ ковариация), \texttt{[3]} & $\chi^2/\text{dof}$ & 1.04 & 1.36 & DSIC не хуже \\
$Om(z)$-диагностика & хронометры + DESI $D_H/r_d$, \texttt{[4]} & знак $\Delta\chi^2$ & --- & --- & \textbf{не разрешает} (зависит от $H_0$) \\
Высокий $z$ ($z > 3$) & предсказание, \texttt{[5]} & $E(z)$, $D_H/r_d$ & расходится & --- & фальсифицируемо \\
Космическое время & $t(z)$ при $z = 5\text{--}12$, \texttt{[6]} & доля от $\Lambda$CDM & $0.68 \to 0.48$ & 1 & \textbf{зона риска ядра} (числом) \\
Опорные данные Planck & $d_{\text{LSS}} = r^*/\theta^*$, \texttt{[7]} & согласованность & --- & --- & сверено (0.00\%) \\
Расстояние до реликта & замыкание по $\theta^*$, \texttt{[8]} & $100\cdot\theta_{\text{model}}$ & 1.04110 & --- & точно (по построению $\mu_d$) \\
Поздняя Вселенная при $\mu_d$ & поправка при $z \le 4.53$, \texttt{[9]} & Мпк & 0.00 & --- & не затронута (тождественно) \\
Степени свободы $\mu_d$ & вырождение $A \leftrightarrow \mu_d$, \texttt{[9]} & dof & 0 & --- & калибровка при форме из P6 \\
Согласованность P6 & галактики JWST vs порог, \texttt{[10]} & --- & --- & --- & флаг: формулировка «средней среды» \\
Рост структур $f\sigma_8$ & Gold-2017 (18 точек), \texttt{[11]} & $\chi^2/\text{dof}$ & 0.70 & $\sim$0.7 & DSIC-фон не хуже \\
Инвариант $S_8$ & вырождение $(\Omega_{m0}, \sigma_8)$, \texttt{[11]} & $1\sigma$ & $0.761 \pm 0.013$ & $0.76\text{--}0.83$ (набл.) & в наблюдаемом диапазоне \\
\bottomrule
\end{tabularx}
\end{center}

Во всех фитах DSIC и $\Lambda$CDM имеют \textbf{одинаковое число свободных параметров}, поэтому $\Delta\chi^2 = \Delta\text{AIC} = \Delta\text{BIC}$.
\section{Связность ядра: машинная проверка тождеств}

Блок \texttt{[0b]} проверяет, что наблюдаемая физика ядра внутренне замкнута --- расстояния, темп и кинематика следуют из одной геометрии без рассогласований:

\begin{center}
\small
\begin{tabularx}{\textwidth}{@{}>{\raggedright\arraybackslash}X l@{}}
\toprule
\textbf{Тождество} & \textbf{Результат} \\
\midrule
$b(\mu)$ строго убывает на $(\mu_0, 1)$ --- инверсия $z \to \mu_e$ корректна & OK \\
$D_H = d\chi/dz$ (сопутствующая $D_H = c/H$) & макс.\ откл. $5.7\cdot10^{-8}$ \\
замкнутая $\chi(\mu)$ = численный $\int c\, dz'/H(z')$ & макс.\ откл. $4.4\cdot10^{-12}$ \\
$q(\mu_{\text{cross}}) = 0$ при точном корне $\mu_{\text{cross}} = \sqrt{3-\sqrt{5}} = 0.874032$ & OK; $z_t = 0.769$ \\
возраст: $t_0 H_0 = \ln(1/\mu_0)\cdot(2-\mu_0^2)/(1-\mu_0^2) = 0.9241$ & $t_0 = 13.41$ Гыр при $H_0 = 67.4$ \\
\bottomrule
\end{tabularx}
\end{center}

Замкнутые формулы статьи и численная геометрия совпадают на уровне машинной точности: темп расширения, расстояния и точка перехода --- одна и та же геометрия, а не независимые подгонки.

\section{Сверхновые: диаграмма расстояний}

\subsection{Якорный фит: калибраторы + $M$ как явный параметр}

Полный каталог Pantheon+SH0ES (1701 SN; в фит идут 1657, включая 77 калибраторов с цефеидными расстояниями как абсолютный якорь). Калибраторы не выбрасываются, а закрепляют абсолютную шкалу; $M$ фитится явно; $H_0$ --- настоящий вывод из якоря. Полная ковариационная матрица STAT+SYS. Три свободных параметра у каждой модели ($\mu_0$/$\Omega_m$, $H_0$, $M$).
\begin{align*}
\text{DSIC}: &\quad \mu_0 = 0.7600 \pm 0.0091,\;\; H_0 = 73.35 \pm 1.01,\;\; M = -19.243 \pm 0.030,\;\; \chi^2/\text{dof} = 0.8783 \\
\Lambda\text{CDM}: &\quad \Omega_m = 0.3308 \pm 0.0180,\;\; H_0 = 73.53 \pm 1.02,\;\; M = -19.244 \pm 0.030,\;\; \chi^2/\text{dof} = 0.8783
\end{align*}
\begin{equation*}
\Delta\chi^2 = \Delta\text{AIC} = \Delta\text{BIC} = +0.05
\end{equation*}
Погрешности --- $1\sigma$ из гессиана $\chi^2$ в минимуме ($C = 2H^{-1}$); параметр модели измерен с точностью $\sim$1.2\%.

На полностью откалиброванных данных DSIC статистически \textbf{неотличима} от $\Lambda$CDM ($|\Delta\chi^2| \ll 2$). Геометрия одного параметра $\mu_0$ воспроизводит ту же зависимость расстояния от красного смещения, что и $\Lambda$CDM с параметром плотности $\Omega_m$, --- без тёмной энергии. $H_0 \approx 73$ наследует известное «напряжение Хаббла» относительно ранней Вселенной ($H_0 \approx 67$) --- одинаково у обеих моделей: DSIC не создаёт нового напряжения и не устраняет существующего.

\begin{quote}
\textbf{Валидация конвейера.} $\Lambda$CDM-ветка того же пайплайна воспроизводит опубликованные значения анализа Pantheon+SH0ES ($\Omega_m \approx 0.33 \pm 0.018$, $H_0 \approx 73.5 \pm 1.0$). Обе модели считаются одним кодом на одних данных; совпадение контрольной ветки с литературой подтверждает корректность сборки.
\end{quote}

\subsection{Устойчивость: STAT-ONLY против STAT+SYS}

Контроль: тот же якорный фит с матрицей только статистических ошибок (без систематик). Если параметры и форма почти не сдвигаются --- результат не держится на систематиках.

\begin{center}
\small
\begin{tabular}{@{}lccc@{}}
\toprule
& $\mu_0$ / $\Omega_m$ & $H_0$ & $\chi^2/\text{dof}$ \\
\midrule
DSIC, STAT+SYS & 0.7600 & 73.35 & 0.8783 \\
$\Lambda$CDM, STAT+SYS & 0.3308 & 73.53 & 0.8783 \\
DSIC, STAT-ONLY & 0.7661 & 73.04 & 0.9242 \\
$\Lambda$CDM, STAT-ONLY & 0.3466 & 73.22 & 0.9244 \\
\bottomrule
\end{tabular}
\end{center}
\begin{equation*}
\Delta\chi^2\;(\text{DSIC} - \Lambda\text{CDM}),\;\text{STAT-ONLY} = -0.39
\end{equation*}
$\mu_0$ сдвигается лишь с 0.7600 до 0.7661 (в пределах $1\sigma$), форма зависимости сохраняется, DSIC остаётся вровень с $\Lambda$CDM. Рост $\chi^2/\text{dof}$ при удалении систематик --- ожидаемый эффект, одинаковый у обеих моделей.
\section{Совместный фит SN + BAO}

Сверхновые Pantheon+SH0ES (1657 SN, якорь) и BAO DESI DR2 (6 бинов, $0.5 < z < 2.33$) с \textbf{полной $12{\times}12$ ковариацией} $D_M$--$D_H$ (учтена корреляция поперечного и радиального масштаба внутри бина). Масштаб звуковой дорожки $r_d$ --- свободный параметр (не сокращается через отношение). Четыре свободных параметра у каждой модели.
\begin{align*}
\text{DSIC}: &\quad \mu_0 = 0.7610,\;\; H_0 = 73.33,\;\; M = -19.243,\;\; r_d = 136.5\;\text{Мпк},\;\; \text{тотал} = 1461.0 \\
\Lambda\text{CDM}: &\quad \Omega_m = 0.3029,\;\; H_0 = 73.75,\;\; M = -19.247,\;\; r_d = 137.1\;\text{Мпк},\;\; \text{тотал} = 1466.0
\end{align*}
\begin{align*}
&\Delta\chi^2\;\text{тотал}\;(\text{DSIC} - \Lambda\text{CDM}) = -5.0 \\
&\text{BAO } \chi^2/\text{dof}\;(\text{dof} = 12-4 = 8):\quad \text{DSIC} = 1.04,\;\; \Lambda\text{CDM} = 1.36
\end{align*}
На совместных данных DSIC описывает геометрию расширения \textbf{не хуже, а формально чуть лучше} $\Lambda$CDM ($\Delta\chi^2 = -5.0$; по шкале Джеффриса это слабое, неубедительное предпочтение, частично отражающее известное натяжение DESI--$\Lambda$CDM). Согласованность по BAO ($\chi^2/\text{dof} \approx 1$) показывает, что та же геометрия $\mu_0$, что описывает сверхновые, одновременно воспроизводит барионный масштаб.

\begin{quote}
\textbf{Замечание о $r_d$.} При совместном фите $r_d \approx 136\text{--}137$ Мпк ниже стандартных $\approx 147$ Мпк; это следствие высокой $H_0 \approx 73$ из SH0ES-якоря (произведение $r_d\cdot H_0$ фиксируется данными). При ранней шкале $H_0 = 67.4$ обратная задача из тех же $D_M/r_d$, $D_H/r_d$ даёт постоянный по $z$ масштаб: $r_d = 147.6$ Мпк по $D_M$ (разброс 1.7) и $r_d = 149.8$ Мпк по $D_H$ (разброс 1.8) --- согласованно между поперечным и радиальным измерением и в пределах $\sim$1.5\% от стандартного значения. Значение $r_d$ смещается вместе с $H_0$ через «напряжение Хаббла» одинаково в обеих моделях.
\end{quote}

\section{$Om(z)$-диагностика --- честный результат}

Диагностика $Om(z) = (E^2 - 1)/((1+z)^3 - 1)$, $E = H/H_0$. У плоской $\Lambda$CDM $Om(z) = \Omega_m = \text{const}$ (горизонталь); DSIC предсказывает U-образную $Om(z)$ с минимумом около $z \approx 2$. $Om(z)$ строится \textbf{из данных} (хронометры $H(z)$ Moresco + DESI $D_H/r_d$) и сравнивается с обеими кривыми. Хронометры подключены в двух режимах ковариации: DIAG (диагональные ошибки) и FULL (полная ковариация Moresco et al.\ 2020 для 15 коррелированных точек; остальные 18 точек компиляции диагональны).

\begin{center}
\small
\begin{tabular}{@{}ccc@{}}
\toprule
$H_0$ & $\Delta\chi^2$ DIAG & $\Delta\chi^2$ FULL \\
\midrule
64 & $-25.0$ & $-24.2$ \\
66 & $-15.6$ & $-15.2$ \\
68 & $-4.6$ & $-4.7$ \\
70 & $+8.3$ & $+7.6$ \\
72 & $+23.1$ & $+21.6$ \\
74 & $+39.9$ & $+37.7$ \\
76 & $+58.9$ & $+55.8$ \\
\bottomrule
\end{tabular}
\end{center}

\textbf{Знак $\Delta\chi^2$ меняется со знаком выбора $H_0$} в обоих режимах (перемена около $H_0 \approx 68.7$): при низкой (Planck-подобной) нормировке данные формально ближе к DSIC, при высокой (SH0ES) --- к $\Lambda$CDM. Результат полностью определяется нормировкой $H_0$, а не формой $Om(z)$; переход DIAG $\to$ FULL меняет $\Delta\chi^2$ лишь на единицы и вердикта не меняет.

\begin{quote}
\textbf{Вывод.} При текущих данных $Om(z)$-тест \textbf{не разрешает} DSIC и $\Lambda$CDM. Заявлять победу DSIC по этой пробе нельзя. Тест станет решающим только при независимо закреплённой $H_0$ и более точных $H(z)$.
\end{quote}

\begin{quote}
\textbf{Оговорка о ковариации.} Полная матрица --- аккуратная реконструкция метода Moresco et al.\ (2020) на официальных входных таблицах (вшиты в скрипт), а не дословный запуск авторского кода. Диагональ совпадает с публикацией точно; недиагональная структура воспроизводит предписанный рецепт коррелированных систематических компонент. На качественный вывод (знак зависит от $H_0$) это не влияет.
\end{quote}
\section{Фальсифицируемые предсказания DSIC}

Там, где данных пока мало, DSIC и $\Lambda$CDM расходятся настолько, что измерения вынесут вердикт. Параметры из совместного фита ($\mu_0 = 0.7610$, $\Omega_m = 0.3029$). Все сигналы этого раздела --- след \textbf{ядра} (Часть I) и от механики второго этажа не зависят: $D_H = c/H$ определяется только $\mu_0$, а поправка $\mu_d$ тождественно нулевая при $z \le 4.53$ (\S8.3).

\subsection{Форма $Om(z)$}

Значения детерминированы параметрами совместного фита:

\begin{center}
\small
\begin{tabular}{@{}ccc@{}}
\toprule
$z$ & $Om_{\text{DSIC}}$ & $Om_{\Lambda\text{CDM}}$ \\
\midrule
0.2 & 0.3566 & 0.3029 \\
0.5 & 0.3297 & 0.3029 \\
1.0 & 0.3064 & 0.3029 \\
2.0 & 0.3032 & 0.3029 \\
3.0 & 0.3266 & 0.3029 \\
5.0 & 0.4043 & 0.3029 \\
\bottomrule
\end{tabular}
\end{center}

DSIC даёт U-образную $Om(z)$ (минимум около $z \approx 2$, подъём к высоким $z$) --- качественная подпись, которой у плоской $\Lambda$CDM нет. Плоская $Om(z)$, измеренная точно, работала бы против DSIC.

\subsection{Радиальный BAO на высоком $z$ (Lyman-$\alpha$)}

\begin{center}
\small
\begin{tabular}{@{}cccc@{}}
\toprule
$z$ & $D_H/r_d$ DSIC & $D_H/r_d$ $\Lambda$CDM & разница \\
\midrule
2.33 & 8.612 & 8.599 & $+0.1\%$ \\
2.5 & 7.978 & 8.013 & $-0.4\%$ \\
3.0 & 6.449 & 6.615 & $-2.5\%$ \\
3.5 & 5.300 & 5.572 & $-4.9\%$ \\
\bottomrule
\end{tabular}
\end{center}

К $z \approx 3.5$ расхождение достигает $\sim$5\%. Точные Lyman-$\alpha$ BAO в этом диапазоне (окно $2.33 < z < 4.53$ целиком ниже порога отлипания) --- прямой тест ядра.

\subsection{Темп расширения за пределами данных}
\begin{align*}
z = 5\phantom{0}: &\quad E_{\text{DSIC}} / E_{\Lambda\text{CDM}} = 1.15 \quad (+15\%) \\
z = 10: &\quad E_{\text{DSIC}} / E_{\Lambda\text{CDM}} = 1.47 \quad (+47\%)
\end{align*}
DSIC предсказывает существенно более жёсткую раннюю историю расширения. Прямое следствие для космического времени --- раздел 7.
\section{Космическое время $t(z)$ --- зона риска ядра, предъявленная числом}

В DSIC собственное время наблюдателя от каши до эпохи $\mu_e$: $t(z)\cdot H_0 = \ln(1/\mu_e(z))\cdot(2-\mu_0^2)/(1-\mu_0^2)$; при $z = 0$ воспроизводит $t_0 H_0 \approx 0.925$ из блока \texttt{[0b]}. Сравнение на общей шкале $H_0 = 67.4$ ($\Omega_m = 0.3029$ из совместного фита, отчего $t_0$ $\Lambda$CDM здесь 13.95 Гыр --- обе модели поставлены в равные условия; вклад излучения $\lesssim$1\% опущен, что сравнение в пользу DSIC не искажает):
\begin{equation*}
\text{Возраст сегодня:}\quad t_{0,\text{DSIC}} = 13.38\;\text{Гыр},\quad t_{0,\Lambda\text{CDM}} = 13.95\;\text{Гыр}
\end{equation*}

\begin{center}
\small
\begin{tabular}{@{}cccc@{}}
\toprule
$z$ & $t_{\text{DSIC}}$, Гыр & $t_{\Lambda\text{CDM}}$, Гыр & отношение \\
\midrule
3.0 & 1.727 & 2.184 & 0.79 \\
5.0 & 0.812 & 1.194 & 0.68 \\
6.0 & 0.604 & 0.948 & 0.64 \\
7.0 & 0.467 & 0.776 & 0.60 \\
7.5 & 0.415 & 0.709 & 0.59 \\
8.0 & 0.371 & 0.651 & 0.57 \\
10.0 & 0.250 & 0.482 & 0.52 \\
12.0 & 0.180 & 0.375 & 0.48 \\
\bottomrule
\end{tabular}
\end{center}

Контекст: квазары с $M_{\text{BH}} \sim 10^9\,M_\odot$ наблюдаются к $z \approx 7.5$; галактики JWST подтверждены до $z \approx 10\text{--}14$; реионизация завершается к $z \approx 6$.

\textbf{Солпитеровская прикидка.} На эддингтоновском пределе e-складывание массы чёрной дыры занимает $\sim$45 млн лет; рост от звёздного зародыша $\sim$$10^2\,M_\odot$ до $\sim$$10^9\,M_\odot$ требует $\sim$16 e-складываний $\approx 0.72$ Гыр. $\Lambda$CDM к $z \approx 7.5$ располагает 0.71 Гыр --- впритык (что уже мотивирует тяжёлые зародыши); DSIC располагает 0.42 Гыр $\approx 9$ e-складываний --- обязательны либо зародыши прямого коллапса $\sim$$10^5\,M_\odot$ при непрерывной аккреции, либо супер-эддингтоновские режимы. Это не опровержение (супер-эддингтоновская аккреция всерьёз обсуждается, в том числе из-за «перемассивных» чёрных дыр JWST), но это \textbf{самое узкое место ядра} --- оптическая поправка Части II его не смягчает, поскольку $H(z)$ не меняет.

\textbf{Оговорка о возрасте сегодня.} Часто цитируемые $13.80 \pm 0.02$ Гыр --- возраст, выведенный внутри $\Lambda$CDM из данных Planck; мерить им не-$\Lambda$CDM модель отчасти циклично. Модельно-независимые звёздные возрасты (старейшие шаровые скопления, $\sim$$13.5 \pm 0.5$ Гыр) значение $t_0 = 13.38$ Гыр проходят.
\section{Часть II: расстояние до реликта (порог отлипания $\mu_d$)}

\begin{quote}
\textbf{Статус.} $\mu_d$ --- константа второго этажа: форма поправки зафиксирована постулатом P6 (амплитуда $A \equiv 1$, резкий порог), значение калибруется по \textbf{одному} наблюдению. Проверка ниже \textbf{не доказывает} $\mu_d$; она подтверждает, что арифметика воспроизводима, опорная цифра Planck верна, поздняя Вселенная не затронута, и честно показывает вырождение, которое фиксирует именно постулат, а не данные.
\end{quote}

\subsection{Опорные данные Planck (самопроверка)}
\begin{align*}
&z^* = 1089.92,\quad r^* = 144.43\;\text{Мпк},\quad 100\theta^* = 1.04110 \\
&d_{\text{LSS}}\;(\text{вшито}) = 13872.8 \pm 25.0\;\text{Мпк} \\
&d_{\text{LSS}} = r^*/\theta^* = 13872.8\;\text{Мпк}\quad (\text{расхождение } 0.00\%)
\end{align*}
Прямо измеренная величина --- угол $\theta^*$. Звуковой горизонт $r^* = 144.43$ Мпк --- величина, \textbf{вычисленная внутри $\Lambda$CDM} из физики фотон-барионной плазмы, которой в DSIC нет; он используется как внешний вход, и это помечено явно. Происхождение акустического масштаба $\sim$147 Мпк --- открытый вопрос ядра.

\subsection{Недобор ядра и калибровка порога по $\theta^*$}

Параметры ядра на ранней шкале --- те же, что для BAO, хронометров и возраста: $\mu_0 = 0.76$, $H_0 = 67.4$, $c/k = 14978.2$ Мпк.
\begin{align*}
&\mu_{\text{cmb}}(z^* = 1089.92) = 0.9999994681 \\
&D_M \text{ база до CMB} = 10936.2\;\text{Мпк} \\
&\text{недобор} = 2936.6\;\text{Мпк}\;\;(21.2\%) \\[4pt]
&\text{калибровка:}\;\; \theta_{\text{model}}(\mu_d) = \theta^* \;\Leftrightarrow\; D_M^{\text{испр}}(\mu_{\text{cmb}}) = r^*/\theta^* = 13872.8\;\text{Мпк} \\
&\mu_d = 0.980640 \\
&z_{\text{detach}} = 4.526 \\
&100\cdot\theta_{\text{model}} = 1.04110 \quad (\text{Planck: } 1.04110)\;\text{--- замыкание точное}
\end{align*}

\subsection{Поздняя Вселенная не затронута; вырождение без P6}

Поправка тождественно нулевая ниже порога:

\begin{center}
\small
\begin{tabular}{@{}ccc@{}}
\toprule
$z$ & $D_M$ база, Мпк & поправка, Мпк \\
\midrule
0.50 & 1932.8 & $+0.0$ \\
1.00 & 3384.2 & $+0.0$ \\
2.33 & 5792.7 & $+0.0$ \\
4.00 & 7335.6 & $+0.0$ \\
4.53 & 7653.3 & $+0.0$ \\
4.63 & 7707.7 & $+49.3$ $\leftarrow$ включена \\
5.00 & 7896.9 & $+220.4$ $\leftarrow$ включена \\
10.00 & 9247.2 & $+1433.1$ $\leftarrow$ включена \\
\bottomrule
\end{tabular}
\end{center}

Все 6 бинов DESI ($z \le 2.33$) поправки не чувствуют --- результаты Части I (разделы 3--6) сохранены полностью.

Без постулата P6 калибровка вырождена --- уравнение $A\cdot(c/k)\cdot[\arcsin\mu_{\text{cmb}} - \arcsin\mu_d] = \text{недобор}$ имеет кривую решений:

\begin{center}
\small
\begin{tabular}{@{}ccc@{}}
\toprule
$A$ & $\mu_d$ & $z_{\text{detach}}$ \\
\midrule
0.5 & 0.923705 & 1.506 \\
0.8 & 0.969868 & 3.344 \\
\textbf{1.0} & \textbf{0.980640} & \textbf{4.526} $\leftarrow$ фиксировано P6 \\
1.2 & 0.986514 & 5.690 \\
2.0 & 0.995097 & 10.266 \\
5.0 & 0.999190 & 26.922 \\
\bottomrule
\end{tabular}
\end{center}

Одно наблюдение фиксирует один параметр: \textbf{степеней свободы ноль}. При форме, зафиксированной P6, $\mu_d$ --- измеренная константа второго этажа (статус как у $\mu_0$ в ядре); предсказанием $z_{\text{detach}}$ не является.

\subsection{След порога, горизонт и досягаемость зондов}

\begin{center}
\small
\begin{tabular}{@{}ccc@{}}
\toprule
$z$ & объект & поправка к $D_M$ \\
\midrule
4.53 & порог отлипания & $+0.0\%$ \\
5.00 & галактики JWST & $+2.8\%$ \\
6.00 & квазары & $+7.1\%$ \\
8.00 & реионизация & $+12.4\%$ \\
11.00 & первые галактики & $+16.6\%$ \\
1089.92 & CMB (калибровка) & $+26.9\%$ \\
\bottomrule
\end{tabular}
\end{center}

Горизонт частиц: база $10953.6$ Мпк; с поправкой $13905.7$ Мпк --- на $32.9$ Мпк дальше поверхности последнего рассеяния: в исправленной геометрии реликт лежит почти вплотную к горизонту.

Досягаемость зондов: DESI DR2 ($z \le 2.33$), DESI DR3 ($\sim$2027, $z \le 3.55$), DESI-II/LBG ($\sim$2029, $2.5 < z < 3.5$) --- все ниже порога. Радиальный BAO в окне $2.33\text{--}4.53$ --- след ядра, не $\mu_d$; сам порог влияет лишь на $D_M$ выше $z \approx 4.5$, где точных линеек, кроме CMB-калибровки, нет. Второй независимый якорь для $\mu_d$ недосягаем до появления точной шкалы расстояний на $z \gtrsim 5\text{--}6$.

\subsection{Флаг согласованности формулировки P6}

Наблюдаемые галактики (JWST, $z \approx 10\text{--}14$; $\mu_e(z{=}14) = 0.9972$) лежат \textbf{выше} порога отлипания ($\mu_d = 0.9806$, $z_{\text{detach}} = 4.53$) --- существуют до отлипания. Буквальное чтение «раздельных объектов не существует при $\mu > \mu_d$» противоречило бы наблюдениям; принятая в статье формулировка P6 относит порог к \textbf{средней среде вдоль пути света}: локальные сгустки с максимальной связанностью падают в себя и коллапсируют в объекты раньше порога, не меняя средней доступности пространства. Скрипт выводит этот флаг явно; арифметика расстояний и результаты Части I от него не зависят.
\section{Часть III: рост структур ($f\sigma_8$)}

Ядро DSIC задаёт только \textbf{фон} --- темп расширения $H(z)$. Отдельный вопрос --- совместим ли этот фон с наблюдаемым \textbf{ростом неоднородностей} (галактик, скоплений). Блок \texttt{[11]} делает минимальный, честно ограниченный шаг: на фоне DSIC решается \textbf{стандартное} уравнение линейного роста возмущений материи в суб-горизонтном GR-пределе,
\begin{equation*}
\frac{d^2\delta}{d(\ln a)^2} + \left[2 + \frac{d\ln H}{d\ln a}\right]\frac{d\delta}{d(\ln a)} - \frac{3}{2}\,\Omega_m(a)\,\delta = 0,
\end{equation*}
и вычисляется наблюдаемая комбинация $f\sigma_8(z) = f\cdot\sigma_8\cdot\delta(z)/\delta(0)$, $f = d\ln\delta/d\ln a$. Сравнение --- с компиляцией \textbf{Gold-2017} (Nesseris, Pantazis \& Perivolaropoulos 2017; 18 слабо коррелированных точек, ковариация трёх точек WiggleZ учтена).

\begin{quote}
\textbf{Статус --- проверка на выживание, не вывод из ядра.} DSIC не выводит уравнение роста из первых принципов: ядро плотностей не вводит, поэтому $\Omega_{m0}$ входит здесь как параметр роста, а само уравнение берётся в GR-пределе на геометрии DSIC. Результат снимает возражение о грубой несовместимости, но не является предсказанием из ядра.
\end{quote}

\subsection{Результат и вырождение $(\Omega_{m0}, \sigma_8)$}

Данные $f\sigma_8$ ограничивают не $\Omega_{m0}$ и $\sigma_8$ по отдельности, а их комбинацию $S_8 = \sigma_8\cdot\sqrt{\Omega_{m0}/0.3}$. Скан вдоль вырожденной долины ($\mu_0 = 0.76$ фиксировано из фона):

\begin{center}
\small
\begin{tabular}{@{}cccc@{}}
\toprule
$\Omega_{m0}$ & $\sigma_8^*$ & $\chi^2/\text{dof}$ & $S_8 = \sigma_8\cdot\sqrt{\Omega_{m0}/0.3}$ \\
\midrule
0.15 & 1.040 & 0.72 & 0.735 \\
0.20 & 0.925 & 0.70 & 0.755 \\
0.25 & 0.835 & 0.73 & 0.762 \\
0.30 & 0.770 & 0.78 & 0.770 \\
0.35 & 0.715 & 0.84 & 0.772 \\
0.40 & 0.670 & 0.91 & 0.774 \\
\bottomrule
\end{tabular}
\end{center}

\begin{align*}
&\text{Формальный минимум:}\;\; \Omega_{m0} = 0.185,\;\; \sigma_8 = 0.955,\;\; \chi^2/\text{dof} = 0.70\;\;(18\text{ точек}) \\
&\text{Инвариант вдоль долины:}\;\; S_8 = 0.761 \pm 0.013 \\
&\text{Срез при Planck-подобном } \Omega_{m0} = 0.30:\;\; \sigma_8 = 0.770,\;\; \chi^2/\text{dof} = 0.73,\;\; S_8 = 0.770
\end{align*}

Вдоль всей долины $\chi^2/\text{dof} \approx 0.7\text{--}0.9$ --- фон DSIC описывает рост структур не хуже $\Lambda$CDM. Инвариант $S_8 = 0.761 \pm 0.013$ устойчив и согласуется с независимыми оценками (Planck $S_8 \approx 0.83$, слабое линзирование $S_8 \approx 0.76$); DSIC-фон попадает в наблюдаемый диапазон без натяжки.

\subsection{Кривая $f\sigma_8(z)$ (срез $\Omega_{m0} = 0.30$)}

Форма $f\sigma_8(z)$ DSIC пологая, с максимумом около $z \approx 0.4\text{--}0.6$, и проходит сквозь облако данных:

\begin{center}
\small
\begin{tabular}{@{}ccc@{}}
\toprule
$z$ & $f\sigma_8$ модель & $f\sigma_8$ набл. $\pm$ ошибка \\
\midrule
0.02 & 0.399 & $0.428 \pm 0.046$ \\
0.15 & 0.421 & $0.490 \pm 0.145$ \\
0.38 & 0.440 & $0.440 \pm 0.060$ \\
0.59 & 0.440 & $0.488 \pm 0.060$ \\
0.73 & 0.434 & $0.437 \pm 0.072$ \\
0.86 & 0.425 & $0.400 \pm 0.110$ \\
1.40 & 0.375 & $0.482 \pm 0.116$ \\
\bottomrule
\end{tabular}
\end{center}

\textit{(таблица прорежена; полный вывод по 18 точкам --- в логе блока \texttt{[11]}).}

\begin{quote}
\textbf{Упрощения (заявлены честно).} (1) Поправка Alcock--Paczynski за разницу фоновых расстояний DSIC vs опорной космологии каждого RSD-измерения не применялась; при $z < 1$ AP-фактор обычно $< 2\text{--}3\%$, мал против ошибок. (2) Уравнение роста --- GR-предельное; вывод роста из связанности (\S12 статьи) остаётся направлением развития. Оба упрощения отмечены как источники систематики, на качественный вывод (совместимость) не влияют.
\end{quote}
\section{Итог}

\noindent\textbf{Что проверено и сходится:}
\begin{itemize}
  \item внутренние тождества ядра выполняются с машинной точностью ($\le 5.7\cdot10^{-8}$);
  \item сверхновые (якорный фит): $\Delta\chi^2 = +0.05$ --- неотличимо от $\Lambda$CDM; параметр измерен: $\mu_0 = 0.7600 \pm 0.0091$;
  \item валидация конвейера: $\Lambda$CDM-ветка воспроизводит опубликованные значения Pantheon+SH0ES;
  \item устойчивость к способу учёта ошибок (STAT-ONLY): $\Delta\chi^2 = -0.39$, сдвиг $\mu_0$ в пределах $1\sigma$;
  \item SN + BAO совместно: $\Delta\chi^2 = -5.0$ (слабое формальное предпочтение DSIC); BAO $\chi^2/\text{dof}$ = 1.04 против 1.36;
  \item линейка $r_d \approx 148$ Мпк при ранней шкале, постоянна по $z$;
  \item второй этаж: замыкание до $\theta^*$ точное, поздняя Вселенная не затронута тождественно, горизонт конечен и лежит в 33 Мпк за реликтом;
  \item рост структур: фон DSIC описывает $f\sigma_8$ (Gold-2017) не хуже $\Lambda$CDM ($\chi^2/\text{dof} \approx 0.7$), инвариант $S_8 = 0.761 \pm 0.013$ в наблюдаемом диапазоне.
\end{itemize}

\noindent\textbf{Что не подтверждено и заявлено честно:}
\begin{itemize}
  \item $Om(z)$-тест при текущих данных не разрешает модели --- знак $\Delta\chi^2$ зависит от $H_0$ (в обоих режимах ковариации);
  \item $\mu_d$ --- калибровка по одному наблюдению (0 степеней свободы) при форме, зафиксированной постулатом; $r^*$ заимствован из $\Lambda$CDM;
  \item формулировка P6 требует чтения «средней среды» (флаг \S8.5);
  \item рост структур --- проверка на выживание, а не вывод из ядра: уравнение роста берётся в GR-пределе, $\Omega_{m0}$ входит как параметр роста (\S9).
\end{itemize}

\noindent\textbf{Где модель встретит решающие проверки:}
\begin{itemize}
  \item форма $Om(z)$; Lyman-$\alpha$ BAO при $z \approx 3\text{--}3.5$ (до $\sim$5\%, досягаемо DESI DR3);
  \item \textbf{космическое время при $z = 5\text{--}12$}: при $z = 7.5$ доступно 0.42 Гыр против 0.71 у $\Lambda$CDM --- солпитеровский бюджет роста чёрных дыр сокращается с $\sim$16 до $\sim$9 e-складываний; это самое узкое место ядра, предъявленное числом.
\end{itemize}

DSIC --- жизнеспособная однопараметрическая геометрическая альтернатива на поздних данных, с чёткими фальсифицируемыми предсказаниями на ранних. Текущие данные её не отвергают и не подтверждают сильнее $\Lambda$CDM; решающими станут высокий $z$, независимо закреплённая $H_0$ и бюджет космического времени.

\section{Воспроизводимость}

Все результаты получены \textbf{единым скриптом} \texttt{dsic\_test.py}; запуск одной ячейкой в Colab либо \texttt{python dsic\_test.py}. Скрипт сам скачивает публичные каталоги и считает $\chi^2$ с полной ковариацией; происхождение каждой вшитой цифры задокументировано в самом файле.

\begin{quote}
\textbf{О времени прогона.} Полный прогон занимает ориентировочно 10--15 минут: скачивание каталога сверхновых ($\sim$65 МБ), фиты Pantheon+SH0ES (несколько минут) и блок \texttt{[11]} (рост структур) --- самый долгий этап (скан по сетке с численным интегрированием ОДУ в каждом узле, $\sim$5 минут). Между печатью блока \texttt{[10]} и результатами \texttt{[11]} возможна длинная тихая пауза --- это ожидаемо, скрипт считает, прерывать не нужно.
\end{quote}

Данные: \textbf{Pantheon+SH0ES} (\texttt{Pantheon+SH0ES.dat}, \texttt{STAT+SYS.cov}, \texttt{STATONLY.cov}) из репозитория \texttt{PantheonPlusSH0ES/DataRelease} (автоскачивание); \textbf{DESI DR2} BAO (Table IV, 6 бинов, вшиты, самопроверка \texttt{[0]}); космические хронометры --- компиляция 33 точек с полной ковариацией Moresco et al.\ (2020) для 15 точек Moresco; \textbf{Planck 2018} ($z^*$, $r^*$, $\theta^*$, $d_{\text{LSS}}$, самопроверка \texttt{[7]}); \textbf{Gold-2017} $f\sigma_8$ (18 точек, Nesseris et al.\ 2017; из Table I Sagredo et al.\ 2018; ковариация WiggleZ --- Blake et al.\ 2012), вшиты в блок \texttt{[11]}.

\begin{center}
\small
\begin{tabularx}{\textwidth}{@{}l >{\raggedright\arraybackslash}X l@{}}
\toprule
\textbf{Блок} & \textbf{Что делает} & \textbf{Раздел} \\
\midrule
\texttt{[0]} & Самопроверка вшитого DESI против Table IV & \S11 \\
\texttt{[0b]} & Машинная проверка внутренних тождеств ядра & \S2 \\
\texttt{[1]} & Якорный SN-фит (STAT+SYS), $M$ явно, $1\sigma$ из гессиана & \S3.1 \\
\texttt{[2]} & Контроль STAT-ONLY & \S3.2 \\
\texttt{[3]} & Совместный SN+BAO, свободный $r_d$, линейка $r_d$ & \S4 \\
\texttt{[4]} & $Om(z)$: устойчивость к $H_0$, ковариации DIAG + FULL & \S5 \\
\texttt{[5]} & Фальсифицируемые предсказания на высоком $z$ & \S6 \\
\texttt{[6]} & Космическое время $t(z)$ DSIC vs $\Lambda$CDM & \S7 \\
\texttt{[7]} & Самопроверка опорных данных Planck & \S8.1 \\
\texttt{[8]} & Недобор ядра, фиксация $\mu_d$ по $\theta^*$ & \S8.2 \\
\texttt{[9]} & Контроль нуля ниже порога; вырождение $A \leftrightarrow \mu_d$ & \S8.3 \\
\texttt{[10]} & След $\mu_d$, горизонт, досягаемость, флаг P6 & \S8.4--8.5 \\
\texttt{[11]} & Рост структур $f\sigma_8$ на фоне DSIC (GR-предел), инвариант $S_8$ & \S9 \\
\bottomrule
\end{tabularx}
\end{center}

Зависимости: \texttt{numpy}, \texttt{scipy} (в Colab предустановлены).

\end{document}
